\documentclass[12pt]{article}
\usepackage{amsmath, amssymb}
\usepackage[margin=1in]{geometry}
\setlength{\parskip}{6pt}
\title{QUBO Formulation: Minimum Dominating Set Problem}
\date{}
\begin{document}
\maketitle

\subsection*{Minimum Dominating Set Problem}

\paragraph{Problem Statement.}
Given an undirected graph $G = (V, E)$, the Minimum Dominating Set problem seeks a subset of vertices $D \subseteq V$ of minimum cardinality such that every vertex in $V$ is either contained in $D$ or adjacent to at least one vertex in $D$. Formally, we require that for every vertex $v \in V$, there exists a vertex $u \in D$ such that $u = v$ or $(u, v) \in E$.

\paragraph{Decision Variables.}
We introduce a binary variable $x_i \in \{0, 1\}$ for each vertex $i \in V$. The variable $x_i$ indicates membership in the dominating set:
\[
x_i = \begin{cases}
1 & \text{if vertex } i \in D, \\
0 & \text{otherwise}.
\end{cases}
\]

\paragraph{QUBO Derivation.}
We construct a Quadratic Unconstrained Binary Optimization (QUBO) formulation by encoding the objective and constraints into a single Hamiltonian $H(\mathbf{x})$.

\textbf{Step 1: Objective Function.}
The goal is to minimize the size of the dominating set, which corresponds to minimizing the sum of the decision variables:
\begin{equation}
\min \sum_{i \in V} x_i.
\end{equation}

\textbf{Step 2: Constraints.}
For each vertex $j \in V$, let $N(j) = \{ i \in V \mid (i, j) \in E \}$ denote the set of neighbors of $j$. The domination requirement is that $j$ must be selected or adjacent to a selected vertex:
\begin{equation}
x_j + \sum_{i \in N(j)} x_i \ge 1, \quad \forall j \in V.
\end{equation}
Define the local domination sum $S_j$ as:
\begin{equation}
S_j = x_j + \sum_{i \in N(j)} x_i.
\end{equation}
The constraint is satisfied if and only if $S_j \ge 1$. Since $x_i$ are binary, $S_j$ is a non-negative integer. A violation occurs only when $S_j = 0$.

\textbf{Step 3: Penalty Terms.}
We enforce the constraints using quadratic penalty terms with weight $\lambda > 0$. For each $j \in V$, we add:
\begin{equation}
P_j(\mathbf{x}) = \lambda \left( 1 - S_j \right)^2.
\end{equation}
This term is zero when $S_j \ge 1$ and equals $\lambda$ when $S_j = 0$. Expanding $P_j$:
\begin{equation}
P_j(\mathbf{x}) = \lambda \left( 1 - 2S_j + S_j^2 \right).
\end{equation}
Using the idempotent property $x_k^2 = x_k$ for binary variables, we expand $S_j^2$:
\begin{align}
S_j^2 &= \left( x_j + \sum_{i \in N(j)} x_i \right)^2 \\
&= x_j^2 + 2x_j \sum_{i \in N(j)} x_i + \left( \sum_{i \in N(j)} x_i \right)^2 \\
&= x_j + 2\sum_{i \in N(j)} x_j x_i + \sum_{i \in N(j)} x_i^2 + \sum_{\substack{i, k \in N(j) \\ i \neq k}} x_i x_k \\
&= x_j + 2\sum_{i \in N(j)} x_j x_i + \sum_{i \in N(j)} x_i + \sum_{\substack{i, k \in N(j) \\ i \neq k}} x_i x_k.
\end{align}
Substituting this back into the penalty expression:
\begin{align}
P_j(\mathbf{x}) &= \lambda \left[ 1 - 2\left( x_j + \sum_{i \in N(j)} x_i \right) + x_j + 2\sum_{i \in N(j)} x_j x_i + \sum_{i \in N(j)} x_i + \sum_{\substack{i, k \in N(j) \\ i \neq k}} x_i x_k \right] \\
&= \lambda \left[ 1 - x_j - \sum_{i \in N(j)} x_i + 2\sum_{i \in N(j)} x_j x_i + \sum_{\substack{i, k \in N(j) \\ i \neq k}} x_i x_k \right].
\end{align}

\textbf{Step 4: Combined Hamiltonian.}
The total QUBO Hamiltonian is the sum of the objective and all penalty terms:
\begin{equation}
H(\mathbf{x}) = \sum_{i \in V} x_i + \sum_{j \in V} P_j(\mathbf{x}).
\end{equation}
Collecting terms, we obtain:
\begin{equation}
H(\mathbf{x}) = c + \sum_{i \in V} Q_{i,i} x_i + \sum_{\substack{i, j \in V \\ i < j}} Q_{i,j} x_i x_j,
\end{equation}
where the constant offset $c$ and coefficients $Q_{i,j}$ are derived below.

\textbf{Step 5: QUBO Coefficients.}
Summing the linear contributions from $P_j$ over all $j \in V$:
\begin{align}
\sum_{j \in V} \left( -\lambda x_j - \lambda \sum_{i \in N(j)} x_i \right) &= -\lambda \sum_{j \in V} x_j - \lambda \sum_{j \in V} \sum_{i \in N(j)} x_i \\
&= -\lambda \sum_{i \in V} x_i - \lambda \sum_{i \in V} x_i \deg(i) \\
&= -\lambda \sum_{i \in V} \left( 1 + \deg(i) \right) x_i,
\end{align}
where $\deg(i) = |N(i)|$ is the degree of vertex $i$. Adding the objective term $\sum_{i \in V} x_i$, the diagonal QUBO entries are:
\begin{equation}
Q_{i,i} = 1 - \lambda \left( 1 + \deg(i) \right), \quad \forall i \in V.
\end{equation}
For off-diagonal entries $i \neq j$, contributions arise from the quadratic terms in $P_j$. The term $x_i x_j$ appears in $P_k$ if $i, j \in N(k)$ or if $j \in N(i)$ or $i \in N(j)$. Specifically:
\begin{itemize}
    \item If $(i, j) \in E$, the term $2\lambda x_i x_j$ appears in $P_i$ and $P_j$, contributing $4\lambda$.
    \item If $(i, j) \notin E$, the term $x_i x_j$ appears in $P_k$ for each common neighbor $k \in N(i) \cap N(j)$. Let $C_{ij} = |N(i) \cap N(j)|$ be the number of common neighbors. Each such $k$ contributes $2\lambda x_i x_j$ via the cross-term in $S_k^2$.
\end{itemize}
Thus, for $i < j$:
\begin{equation}
Q_{i,j} = \begin{cases}
4\lambda & \text{if } (i, j) \in E, \\
2\lambda \, C_{ij} & \text{if } (i, j) \notin E.
\end{cases}
\end{equation}
The constant offset is the sum of constants from all penalties:
\begin{equation}
c = \sum_{j \in V} \lambda = \lambda |V|.
\end{equation}

\paragraph{Penalty Weight Justification.}
The penalty weight $\lambda$ must be chosen large enough to ensure that any infeasible solution has a strictly higher energy than the optimal feasible solution. The maximum possible value of the objective function is $|V|$. Violating a single constraint increases the penalty by at least $\lambda$, while the maximum reduction in the objective from violating constraints is bounded by $|V|$. Therefore, choosing $\lambda > |V|$ guarantees that the penalty for any infeasible assignment dominates any potential gain in the objective, ensuring the global minimum of $H(\mathbf{x})$ corresponds to a feasible and optimal dominating set.

\paragraph{Correctness Argument.}
Minimizing $H(\mathbf{x})$ over $\mathbf{x} \in \{0, 1\}^{|V|}$ recovers the optimal dominating set because:
\begin{enumerate}
    \item \textit{Feasibility:} For any feasible solution, all constraints are satisfied, so all penalty terms vanish. The energy equals the objective value.
    \item \textit{Optimality:} For any infeasible solution, at least one constraint is violated, incurring a penalty of at least $\lambda$. Since $\lambda > |V| \ge \text{objective value}$ for any solution, the energy of any infeasible solution exceeds the energy of any feasible solution. Thus, the global minimum must be feasible. Among feasible solutions, $H(\mathbf{x})$ reduces to the objective function, so the minimum corresponds to the minimum dominating set.
\end{enumerate}

\paragraph{Illustrative Example.}
Consider a path graph with $N=3$ vertices $V=\{0, 1, 2\}$ and edges $E=\{(0, 1), (1, 2)\}$. The degrees are $\deg(0)=1$, $\deg(1)=2$, $\deg(2)=1$. The number of common neighbors is $C_{02}=1$ (vertex $1$), and $C_{01}=C_{12}=0$. We choose $\lambda = 10$, satisfying $\lambda > N$.

Using the general formulas:
\begin{align}
Q_{0,0} &= 1 - 10(1 + 1) = -19, \\
Q_{1,1} &= 1 - 10(1 + 2) = -29, \\
Q_{2,2} &= 1 - 10(1 + 1) = -19, \\
Q_{0,1} &= 4 \times 10 = 40, \\
Q_{1,2} &= 4 \times 10 = 40, \\
Q_{0,2} &= 2 \times 10 \times 1 = 20, \\
c &= 10 \times 3 = 30.
\end{align}
The resulting QUBO Hamiltonian is:
\begin{equation}
H(x_0, x_1, x_2) = 30 - 19x_0 - 29x_1 - 19x_2 + 40x_0x_1 + 40x_1x_2 + 20x_0x_2.
\end{equation}
Evaluating $H$ for candidate solutions:
\begin{itemize}
    \item $\mathbf{x} = (0, 1, 0)$: $H = 30 - 29 = 1$. This corresponds to dominating set $\{1\}$, size $1$.
    \item $\mathbf{x} = (1, 0, 1)$: $H = 30 - 19 - 19 + 20 = 12$. This corresponds to dominating set $\{0, 2\}$, size $2$.
    \item $\mathbf{x} = (0, 0, 0)$: $H = 30$. Infeasible.
\end{itemize}
The minimum energy is $1$, achieved at $\mathbf{x} = (0, 1, 0)$, correctly identifying the minimum dominating set $\{1\}$.

\end{document}